\scanpage{25}
\uline{Corollaire} 1 Si $v\in D_n^{\deleted{\unsure{*}}}\smallsetminus D_n^0$, alors $\varphi(v)\ne\mathrm{id}_E$

En effet, \deleted{\unsure{$\varphi(v)$}} il est clair que $\varphi(v)$ échange les deux ordres circulaires \inserted{\text{$\{v_E\mid v\in\unsure{w|\omega|u}\}$}} sur $S$, \deleted{distincts des \unsure{propriétés postulées}} \deleted{\unsure{et|donc} $v_E$} donc ne peut être l'identité sur $S$, donc pas non plus sur $E$.

\uline{Corollaire 2} Considérons \deleted{\unsure{le|la}} l'entier $n'$ comme \inserted{\text{dans (b)}} dessus. Alors \deleted{$\varphi':D_{n'}$}
\[
\left\{\begin{aligned}
\varphi' &:D_{n'}\longrightarrow\mathrm{Aff}(E)\\
s_0&\in E
\end{aligned}\right.
\]
définissent dans $E$ un polygône épinglé à $n'$ côtés.

\uline{Scholie} On peut considérer que la donnée $\{\varphi,s_0\}$ correspond à une réalisation géom du polygone comb. type $\Pi_0(n)$, qui est un \uline{revêtement à $d$ feuillets} de la réalisation géom \inserted{\text{$(\varphi',s_0)$}} du polyg. comb. type $\Pi_0(n')$. Quand on fera des calculs avec un
% archive sheet number: 24
